\subsection{GLM-5.2: critic-based RL and parallel OPD for long-horizon tasks
\citep{glm52team2026}}
\paragraph{Source type.} First-party research article.
\paragraph{Research question.} Ask how a post-training system should handle compacted,
variable-length long-horizon trajectories and combine specialized expert behaviors.
\paragraph{Method.} Document the move from group-relative optimization to critic-based PPO
so individual rollouts can train without a fixed same-prompt sibling group. Record its
token-level critic advantages, length-balanced loss over compacted sub-traces, online
anti-hacking, and parallel on-policy distillation from multiple expert models.
\paragraph{Main evidence.} Record benchmark results, efficiency claims, critic or OPD
ablations, and the stated two-day multi-expert OPD run as distinct categories of evidence.
\paragraph{Connection to the position.} Provide a production-scale case in which a critic
removes GRPO's group-sampling requirement and supplies denser credit for ragged long-horizon
traces; this motivation is separate from off-policy trace reconstruction.
\paragraph{Critical assessment.} Treat first-party benchmark claims cautiously; examine
critic-target construction, compaction boundaries, anti-hack false positives, expert conflicts,
and the absence or presence of controlled comparisons.
\paragraph{Takeaway.} Complete after reading.
